The work demonstrated here would have been impossible, and vastly less enjoyable without the guidance, support, and friendship of so many people.

First, I would like to thank my advisor Chris Ballance, who gave me the opportunity to explore this work. He has a knack for pinning down faulty hardware when given a list of bizarre, seemingly-unrelated symptoms --- this has saved me many a frustrated hour in the lab. Aside from this technical knowledge, Chris has also provided the necessary pragmatism in how to spend one's limited time and effort --- often with his first question on proposed lab improvements: ``Do we care?'' This framework for thinking, getting straight to the heart of what we want, has been invaluable. 

I want to thank Raghu Srinivas, who although only this last year has been officially recognised as my co-advisor, has been a constant guide for me throughout my DPhil. He is astoundingly creative in his physics, providing many of the background ideas in this work. I am so proud of his recent Professorship, a well-deserved appointment, and one I am sure many future students will be appreciative of.

As for the \emph{FastGates} team, it is difficult to express my gratitude for working along side its members. When we were ourselves a ``hybrid'' two-apparatus group --- \emph{FastBlades} --- Oana Băzăvan and Sebastian Saner set the dynamic of our group, of which I am so thankful. They provided both a welcoming, fun environment, but also great inspiration for the physicist I wanted to be. Due to both of their creativity and dedication, the \emph{Blade} experiment, although perhaps a bit dusty, was truly demonstrating amazing things. I am grateful that I was able to be a small part of that work. 

I want to thank Sebastian and Mariella for the \emph{FastGates} system design, especially that of the vacuum hardware. I have fond memories of overheating in the cleanroom with Fabian and Sebastian constructing that system; it was certainly stressful, and sometimes I am amazed it all worked out. I would not be surprised if a few molecules of river water are still somewhere in the chamber...

Iver {\O}vergaard, from summer placement, to master student, to DPhil, has had perhaps the biggest impact on my time on the project. I have thoroughly enjoyed the many hours working with him on the apparatus and could not have asked for a better labmate to build this experiment with. 
Many of the results in this thesis were enabled directly by him, and I am looking forward to see what incredible things he does for the remainder of his DPhil.

I want to thank the other current members of the group: Kai Shinbrough, Giorgio Canalella, and Emily Hirsch. Thanks to you all, I am very much excited to get this upcoming vacation over with and join you back in the lab.

It takes a large support system to build-up a lab. I am thankful of the entire Oxford Ion Trapping Group, captained by David Lucas. I am particularly appreciative of Gabriel Araneda for his extensive advice and help over the years, and to Lorenzo Versini for his work on the shared Beecroft laser systems. Thanks to David Nadlinger, Peter Drmota, Andres Vazquez Brennan, Jamie Leppard, and Tim Wohlers-Reichel for an assortment of useful chats, technical support and infrastructure. Thank you to Tenzan, the group's resident theorist, for our collaboration on the virtual distillation work.

Finally, I would like to thank the important people outside of the lab. Thank you to my parents; my mum whose curiosity about the world directly shaped my own, and to my dad who I am sure initiated my interest in science at a young age with a set of wooden counting blocks, and his own technical mind. I have been lucky enough to live near family throughout my time in Oxford. Thank you to Claudia, Blake, Casey, Tiffany and Martin, who, although this period of our lives has been busy, we have always found time to spend together. 

To my fianc\'ee Lucy, who I was fortunate enough to meet early on in my time here, and has been a constant source of joy in the last four years. This would not have been possible without you, and I am excited to join your family in just a few short months.\\

Thank you to Raghu, Iver, Kai, Emily, Giorgio, Sebastian, and Oana for proofreading this manuscript.

